%%%%%%%%%%%%%%%%%%%%%%%%%%%%%%%%%%%%%%%%%%%%%%%%%%%%%%%%%%%%%%%%%%%%%%%%%%%%
% CAPÍTULO 14: RECURSOS NECESARIOS
%%%%%%%%%%%%%%%%%%%%%%%%%%%%%%%%%%%%%%%%%%%%%%%%%%%%%%%%%%%%%%%%%%%%%%%%%%%%

\chapter{Recursos necesarios}
\label{ch:recursos}

\section{Recursos software}

Los recursos software previstos se resumen en la Tabla~\ref{tab:recursos-software}.

\begin{table}[H]
  \centering
  \begin{tabularx}{\textwidth}{Y{0.9091} Y{1.3864} Y{0.7045}}
    \hline
    \textbf{Recurso} & \textbf{Uso} & \textbf{Coste} \\
    \hline
    Angular + Ionic + Capacitor & \textit{Frontend} web y móvil con una única base de código & Open Source \\
    Node.js + NestJS & \textit{Backend} y APIs & Open Source \\
    PostgreSQL + PostGIS & Persistencia y consultas geoespaciales & Open Source \\
    GitHub + GitHub Actions & Control de versiones y CI/CD & Gratuito \\
    SonarCloud, CodeQL y Dependabot & Análisis de calidad y seguridad del código en cada cambio & Gratuito \\
    Docker + Docker Compose & Entorno de ejecución reproducible (base de datos, OSRM y servidor de correo de pruebas) & Open Source \\
    OpenStreetMap + OSRM + Leaflet & Visualización de mapas, rutas y desvíos & Open Source \\
    CartoCiudad (IGN) & Búsqueda de direcciones y dirección de un punto & Gratuito (cita obligatoria) \\
    Motor de simulación propio & Generación y ejecución de escenarios & Desarrollo propio \\
    Herramientas de \textit{testing} & Jest + Supertest, Vitest, Cypress, Artillery/k6 & Open Source \\
    TikZ (\LaTeX) & Diagramas UML y de arquitectura & Open Source \\
    Figma (con Figma AI) & Maquetas de las pantallas e ilustración de acceso & Gratuito \\
    \hline
  \end{tabularx}
  \caption{Recursos software.}
  \label{tab:recursos-software}
\end{table}

\textit{Uso de inteligencia artificial generativa:} la ilustración de las pantallas de acceso se generó con Figma AI, la función de inteligencia artificial de Figma, y después se revisó y adaptó a la identidad visual de la plataforma. Las maquetas de las pantallas se diseñaron en Figma.

\section{Recursos hardware}

\begin{table}[H]
  \centering
  \begin{tabularx}{\textwidth}{Y{0.7143} Y{1.2857}}
    \hline
    \textbf{Recurso} & \textbf{Uso} \\
    \hline
    Ordenador de desarrollo & Programación, pruebas y ejecución de simulaciones \\
    Entorno de despliegue local o \textit{cloud} & Pruebas de integración y CI/CD \\
    Dispositivo móvil (opcional) & Pruebas de la aplicación Ionic \\
    \hline
  \end{tabularx}
  \caption{Recursos hardware.}
  \label{tab:recursos-hardware}
\end{table}

\section{Recursos de datos}

\begin{table}[H]
  \centering
  \begin{tabularx}{\textwidth}{Y{0.5714} Y{1.4286}}
    \hline
    \textbf{Recurso} & \textbf{Descripción} \\
    \hline
    Datos sintéticos & Usuarios, viajes, demanda y horarios generados para simulación \\
    Datos geográficos & OpenStreetMap para la red vial y CartoCiudad (IGN) para las direcciones del entorno de la UAL \\
    Datos de la UAL & No disponibles inicialmente; se usarán supuestos documentados \\
    Escenarios reproducibles & Configuraciones y semillas registradas para garantizar la repetibilidad \\
    \hline
  \end{tabularx}
  \caption{Recursos de datos.}
  \label{tab:recursos-datos}
\end{table}

\section{Recursos humanos}

\begin{table}[H]
  \centering
  \begin{tabularx}{\textwidth}{Y{0.7143} Y{1.2857}}
    \hline
    \textbf{Recurso} & \textbf{Rol} \\
    \hline
    Autor del TFM & Diseño, implementación, evaluación y redacción \\
    Director del TFM & Supervisión metodológica y técnica \\
    Usuarios potenciales (opcional) & Pruebas de usabilidad, si el tiempo lo permite \\
    \hline
  \end{tabularx}
  \caption{Recursos humanos.}
  \label{tab:recursos-humanos}
\end{table}

\section{Recursos de conocimiento}

Se prevé utilizar los siguientes recursos de conocimiento:
\begin{itemize}
    \item Literatura del corpus principal (77 publicaciones).
    \item Referencias secundarias (15).
    \item Documentación técnica de las herramientas utilizadas.
    \item Plantilla ERS de la UAL para la especificación de requisitos.
\end{itemize}

\section{Estimación de costes}

El TFM se apoya íntegramente en herramientas \textit{Open Source} y planes gratuitos, por lo que el coste económico directo es prácticamente nulo. El principal recurso es el tiempo dedicado al desarrollo y a la redacción, junto con la supervisión del director.
